\chapter{Metodologia}
\label{sec:metodologia}

Este capítulo descreveu a metodologia da pesquisa. A Seção~\ref{sec:tipopesquisa} classificou a pesquisa quanto à natureza, aos objetivos e ao procedimento técnico, e a Seção~\ref{sec:procedimentos} descreveu as etapas do trabalho.

%-----------------------------------
\subsection{Tipo de Pesquisa}
\label{sec:tipopesquisa}

Quanto à natureza, esta pesquisa foi aplicada. Sua finalidade foi produzir um artefato de software para um problema concreto identificado num contexto institucional real, não ampliar o conhecimento teórico de forma genérica. \citeonline{gil2008} distinguiu a pesquisa aplicada da pesquisa básica precisamente por esse critério: a primeira foi motivada pela necessidade de resolver problemas práticos e seus resultados tiveram aplicação imediata.

Quanto aos objetivos, a pesquisa foi descritiva. O ponto de partida foi a caracterização do processo de agendamento de espaços físicos no IFNMG -- Campus Salinas: como ele funcionava, onde falhava e o que os usuários envolvidos precisavam que mudasse. \citeonline{gil2008} definiu a pesquisa descritiva como aquela que buscava descrever as características de determinado fenômeno e estabelecer relações entre suas variáveis, sem que o pesquisador interferisse nos fatos observados.

Quanto ao procedimento técnico, a pesquisa adota o estudo de caso em combinação com pesquisa bibliográfica. O estudo de caso é adequado porque a investigação recai sobre um ambiente específico, o Campus Salinas, com histórico, atores e restrições próprios que precisam ser compreendidos antes de qualquer proposta de solução. \citeonline{yin2015} define o estudo de caso como uma investigação empírica que examina um fenômeno contemporâneo em profundidade e em seu contexto real, sendo especialmente indicado quando os limites entre o fenômeno e o contexto não são claramente separáveis. A pesquisa bibliográfica, por sua vez, dá sustentação às decisões de projeto, de modelagem e de escolha tecnológica, com base no que a literatura já consolidou sobre engenharia de software, controle de acesso e desenvolvimento web.

%-----------------------------------
\subsection{Procedimentos Metodológicos}
\label{sec:procedimentos}

Esta pesquisa foi organizada em cinco etapas: diagnóstico do problema, levantamento e análise de requisitos, modelagem do sistema, implementação e validação. As etapas foram conduzidas de forma iterativa: a cada reunião com a Direção de Ensino, os requisitos eram refinados e incorporados ao sistema em desenvolvimento. Ao final do trabalho, foram concluídas as cinco etapas: diagnóstico do problema, levantamento e análise de requisitos, modelagem, implementação e validação técnica por meio de testes automatizados.

%-----------------------------------
\subsubsection{Diagnóstico do Problema}
\label{sec:coleta}

A primeira etapa consistiu no diagnóstico do problema. Por meio de conversas semiestruturadas, interações técnicas informais e validações contínuas com a Diretora de Ensino Técnico, responsável pela gestão dos espaços físicos no campus, foram mapeados o processo atual de agendamento, suas limitações operacionais e as necessidades concretas dos usuários envolvidos, como docentes e equipe administrativa, além da consulta pública, disponível sem autenticação. As informações obtidas nessa etapa fundamentaram a especificação dos requisitos apresentada a seguir.

Os dados desta pesquisa são primários e foram coletados diretamente com a responsável pelo processo. O levantamento ocorreu de forma iterativa, por meio de conversas semiestruturadas, interações técnicas, trocas de mensagens e discussões sobre o agendamento. \citeonline{gil2008} aponta a entrevista como uma das técnicas de coleta de dados mais utilizadas nas ciências sociais por permitir que o entrevistador obtenha informações sobre fatos, opiniões, percepções e necessidades de forma direta, com abertura para explorar aspectos não previstos no roteiro inicial. Neste trabalho, essa característica foi estendida para um formato iterativo, em que o diálogo com a fonte se prolongou para além de sessões pontuais.

As interações foram conduzidas com a Diretora de Ensino Técnico do IFNMG -- Campus Salinas, responsável pela gestão dos agendamentos de espaços físicos no campus. Na fase de especificação, foram realizadas três sessões remotas: em julho de 2025, novembro de 2025 e abril de 2026.

A primeira sessão mapeou o processo atual de agendamento, incluindo a forma de solicitar reservas, os responsáveis pelo controle, os conflitos mais frequentes e as informações consultadas pelos usuários. A segunda foi dedicada à validação dos requisitos preliminares. As informações levantadas foram apresentadas de forma organizada para que a Diretora confirmasse, corrigisse ou complementasse o entendimento do pesquisador. A terceira sessão confirmou os requisitos finais e esclareceu questões surgidas durante a modelagem.

Entre essas sessões, houve comunicação técnica por mensagens e reuniões breves. Esses contatos foram usados para esclarecer dúvidas sobre o processo e validar decisões de modelagem à medida que surgiam.

Durante o desenvolvimento, foram realizadas mais três reuniões remotas. Em agosto de 2026, a planilha de horários do campus foi apresentada e os requisitos foram refinados. Em setembro de 2026, uma reunião tratou de novos refinamentos, e outra apresentou o sistema em funcionamento à Direção de Ensino.

A opção pelo formato semiestruturado permitiu seguir um roteiro comum sem limitar a discussão aos temas previstos inicialmente. O caráter iterativo foi necessário porque novas dúvidas surgiram durante a modelagem e precisaram ser confirmadas com a responsável pelo processo. Por isso, o levantamento não se limitou a uma única entrevista ou a um questionário fechado.

Como a Diretora de Ensino era a responsável pelo processo de agendamento e apresentou a necessidade que originou o sistema, ela foi adotada como principal referência institucional para o levantamento e a validação dos requisitos. Os requisitos foram discutidos, refinados e validados nas reuniões realizadas ao longo do desenvolvimento. As percepções dos docentes e dos demais usuários não foram submetidas a uma avaliação formal de usabilidade, o que constituiu uma limitação deste trabalho.

O diagnóstico mostrou que o sistema em uso exige o cadastro manual de todos os horários ocupados antes de qualquer reserva, e esse cadastro é refeito a cada semestre. Os docentes não consultam a disponibilidade nem solicitam reservas sem passar por um servidor, o que concentra todo o controle em um único ponto. Também não há registro consolidado do uso dos espaços que apoie o planejamento da administração.

Dessas limitações derivam as necessidades que orientaram os requisitos: importar a planilha de horários já produzida pelo campus, marcando automaticamente como ocupados os horários das aulas; manter essa ocupação atualizada por reimportação ou ajuste pontual, sem integração em tempo real com sistemas externos; notificar automaticamente os envolvidos no fluxo de reserva; permitir a consulta pública da ocupação; e gerar um mapa da utilização dos espaços ao longo do semestre.

%-----------------------------------
\subsubsection{Levantamento e Análise de Requisitos}
\label{sec:requisitos}

A segunda etapa foi o levantamento e a análise de requisitos. \citeonline{pressman2016} descrevem a engenharia de requisitos como um processo contínuo de identificação, análise, especificação e validação das necessidades dos envolvidos, e ressaltam que a qualidade dessa etapa determina a qualidade do que será desenvolvido depois. A especificação de requisitos foi elaborada a partir das informações coletadas nas reuniões e do referencial bibliográfico consultado, e organizada em duas categorias: requisitos funcionais e requisitos não funcionais.

Os requisitos funcionais descrevem o que o sistema deve fazer, ou seja, as funções e os serviços que ele oferece aos usuários \cite{sommerville2018}. Como o acesso é definido por perfis configuráveis (Seção~\ref{sec:rbac}), cada funcionalidade fica disponível a quem tiver a permissão correspondente, e não a um perfil fixo. A Tabela~\ref{tab:rf} apresenta os requisitos funcionais, agrupados nos módulos: ensalamento (RF01, RF02, RF13--RF16); espaços, cursos, disciplinas e turmas (RF03, RF17); usuários e perfis (RF04, RF18--RF20); agenda de reservas (RF05--RF08, RF21--RF23); notificações e comunicados (RF09, RF10, RF24); consulta pública (RF11, RF28); painel, relatórios e auditoria (RF12, RF25, RF26); configurações (RF27); e central de ajuda (RF29).

\begin{SingleSpace}
\begin{longtable}{@{}p{1.8cm}p{12cm}@{}}
	\caption{Requisitos funcionais do sistema}\label{tab:rf} \\
	\hline
	\textbf{Código} & \textbf{Descrição} \\ \hline
	\endfirsthead
	\multicolumn{2}{@{}l}{\small\textit{Continuação da Tabela~\thetable}} \\ \hline
	\textbf{Código} & \textbf{Descrição} \\ \hline
	\endhead
	\hline
	\endfoot
	RF01 & Importar a planilha de horários institucional do semestre, identificando automaticamente os espaços e os horários ocupados pelas aulas, com prévia do resultado e confirmação antes da gravação \\
	RF02 & Atualizar os horários ocupados ao longo do semestre, por reimportação da planilha ou pela manutenção das aulas no próprio sistema (RF13 a RF16) \\
	RF03 & Cadastrar e manter os espaços físicos do campus, incluindo os outros nomes pelos quais aparecem na planilha e a desativação de espaços fora de uso \\
	RF04 & Gerenciar usuários e atribuir a cada um perfil de acesso \\
	RF05 & Consultar a disponibilidade dos espaços físicos \\
	RF06 & Submeter solicitações de reserva de espaços físicos \\
	RF07 & Aprovar ou rejeitar solicitações de reserva \\
	RF08 & Cancelar reservas já aprovadas \\
	RF09 & Notificar automaticamente, no próprio sistema, os responsáveis pela análise quando uma nova solicitação de reserva for submetida \\
	RF10 & Notificar automaticamente, no próprio sistema, o solicitante sobre o resultado da análise da sua solicitação \\
	RF11 & Disponibilizar consulta pública, sem autenticação, aos horários dos espaços por sala, docente, curso ou turma, com o horário original e o calendário por data, sem expor e-mails nem a descrição das reservas \\
	RF12 & Gerar relatórios sobre a utilização dos espaços, incluindo o mapa de ocupação por dia e horário e o relatório das solicitações de reserva \\
	RF13 & Cadastrar aula recorrente, informando turma, disciplina, docentes, espaço, período de vigência, horário e duração \\
	RF14 & Alterar a recorrência de uma aula (dia, horário, duração, espaço ou docentes) a partir de uma data, sem alterar as ocorrências anteriores \\
	RF15 & Encerrar a recorrência de uma aula a partir de uma data, preservando o histórico das ocorrências anteriores \\
	RF16 & Cancelar a ocorrência de uma aula em uma data, sem alterar a recorrência, pelo administrador ou pelo docente da aula \\
	RF17 & Cadastrar e manter cursos, disciplinas e turmas \\
	RF18 & Criar perfis de acesso e definir as permissões de cada um \\
	RF19 & Importar usuários em lote por meio de arquivo CSV \\
	RF20 & Enviar ao novo usuário, por e-mail, o link para definir a senha, e permitir a recuperação de senha \\
	RF21 & Sugerir ao solicitante uma alternativa à reserva solicitada (outro espaço, data ou horário), que ele pode aceitar ou recusar \\
	RF22 & Registrar solicitação de reserva em nome de outro usuário \\
	RF23 & Acompanhar as próprias solicitações e ver a própria agenda de aulas e reservas \\
	RF24 & Publicar comunicados para todos os usuários, para os usuários de um perfil ou para usuários escolhidos \\
	RF25 & Exibir painel administrativo com as solicitações pendentes, a ocupação dos espaços e as pendências da grade (aulas sem espaço ou sem docente, conflitos de docente, choques entre aulas e reservas) \\
	RF26 & Registrar em auditoria as ações realizadas no sistema e permitir consultá-las com filtros \\
	RF27 & Configurar o servidor de envio de e-mails pela interface do sistema \\
	RF28 & Exportar o calendário público do semestre, com os filtros escolhidos, em arquivo iCalendar (.ics), para importação em aplicativos de agenda \\
	RF29 & Disponibilizar uma central de ajuda com tutoriais passo a passo de cada funcionalidade, com capturas das telas no computador e no celular, mostrando ao usuário só as funcionalidades liberadas para o seu perfil e abrindo direto na ajuda da tela de onde foi acessada \\
\end{longtable}
\end{SingleSpace}

Os requisitos não funcionais estabelecem as restrições de qualidade que o sistema precisa satisfazer \cite{sommerville2018}. Eles foram definidos a partir dos atributos de qualidade descritos pelo autor: confiabilidade e segurança (RNF01, RNF02, RNF07, RNF08, RNF09); eficiência (RNF03, RNF10, RNF11); aceitabilidade (RNF04, RNF05, RNF13); e manutenibilidade (RNF06, RNF12). \citeonline{pressman2016} reforçam que esses requisitos condicionam a arquitetura e devem ser considerados desde o início do projeto. A Tabela~\ref{tab:rnf} apresenta os requisitos não funcionais.

\begin{table}[!htb]
	\centering
	\caption{Requisitos não funcionais do sistema}
	\begin{tabular}{@{}p{1.8cm}p{12cm}@{}}
	\hline
		\textbf{Código} & \textbf{Descrição} \\ \hline
		RNF01 & Controlar o acesso com base em papéis (RBAC), restringindo cada funcionalidade ao perfil correspondente \\
		RNF02 & Proteger os dados dos usuários contra acessos indevidos: proteções contra CSRF, XSS e injeção de SQL, senha do servidor de e-mail armazenada cifrada, recuperação de senha que não revela quais e-mails estão cadastrados e encerramento das sessões do usuário quando outra pessoa altera o perfil ou o e-mail dele \\
		RNF03 & Ter tempo de resposta adequado nas consultas de disponibilidade, mesmo nos períodos de pico \\
		RNF04 & Funcionar no navegador web, sem instalar nada nos dispositivos dos usuários \\
		RNF05 & Ter usabilidade compatível com o perfil tecnológico dos docentes e consistência visual entre as telas, sem exigir treinamento extenso \\
		RNF06 & Ser manutenível, com separação clara entre as camadas de dados, lógica e apresentação \\
		RNF07 & Funcionar sem depender da infraestrutura de horários do campus, sem consultas em tempo real a sistemas externos \\
		RNF08 & Impedir reservas sobrepostas de um mesmo espaço, mesmo quando solicitações são analisadas ao mesmo tempo, por restrição de integridade no banco \\
		RNF09 & Não permitir alterar nem excluir os registros de auditoria \\
		RNF10 & Enviar os e-mails em segundo plano, com novas tentativas em caso de falha, sem atrasar a resposta ao usuário \\
		RNF11 & Exibir as notificações em tempo real, sem recarregar a página \\
		RNF12 & Ser distribuído em contêineres, para facilitar a execução em ambientes diferentes \\
		RNF13 & Adaptar as telas a computadores e celulares \\ \hline
	\end{tabular}
	\label{tab:rnf}
\end{table}

%-----------------------------------
\subsubsection{Modelagem do Sistema}
\label{sec:modelagem}

A modelagem utilizou três diagramas da UML: casos de uso, classes e sequência. Os diagramas refletem a versão implementada do sistema.

O diagrama de casos de uso (Figura~\ref{fig:casos_uso}) organiza as funcionalidades em módulos. Como os perfis do sistema são configuráveis (Seção~\ref{sec:rbac}), cada ator representa um conjunto de permissões, e não um perfil fixo: Solicitante, Analista de reservas, Gestor do ensalamento e Administrador especializam o Usuário autenticado, e o Superadministrador especializa o Administrador. O Visitante representa a consulta pública, sem autenticação.

\begin{figure}[!htb]
    \centering
    \caption{Diagrama de Casos de Uso do Sistema de Agendamento}
    \resizebox{\textwidth}{!}{% Diagrama de casos de uso do sistema de agendamento (gerado por gen_uc.py).
% Os atores representam conjuntos de permissões, e não perfis fixos.
\begin{tikzpicture}[
    font=\footnotesize,
    uc/.style={ellipse, draw, fill=white, minimum width=4.4cm, minimum height=0.8cm,
               align=center, inner xsep=-2pt, inner ysep=2.5pt},
    pkg/.style={draw, rounded corners=2pt, fill=black!4},
    pkglabel/.style={font=\scriptsize\bfseries, anchor=north west, inner sep=3pt},
    rel/.style={dashed, -{Stealth[length=2mm]}},
    rellabel/.style={font=\scriptsize\itshape, right, inner sep=2pt},
    gen/.style={-{Triangle[open, length=2.5mm, width=2.5mm]}},
    lbl/.style={anchor=north, align=center},
    hit/.style={minimum width=0.7cm, minimum height=1.3cm},
    actor/.pic={
        \draw (0,0.55) circle (0.17);
        \draw (0,0.38) -- (0,-0.1) (-0.3,0.22) -- (0.3,0.22)
              (0,-0.1) -- (-0.25,-0.5) (0,-0.1) -- (0.25,-0.5);
    },
]
\draw[thick, rounded corners=4pt] (2.6,-0.55) rectangle (13.4,20.45);
\node[font=\small\bfseries] at (8,20.05) {Sistema de Agendamento de Espaços Físicos};
\draw[pkg] (2.8499999999999996,16.0) rectangle (7.85,19.6);
\node[pkglabel] at (2.8499999999999996,19.6) {Consulta pública};
\draw[pkg] (2.8499999999999996,6.35) rectangle (7.85,15.7);
\node[pkglabel] at (2.8499999999999996,15.7) {Agenda de reservas};
\draw[pkg] (2.8499999999999996,2.45) rectangle (7.85,6.05);
\node[pkglabel] at (2.8499999999999996,6.05) {Comunicação e ajuda};
\draw[pkg] (8.15,12.55) rectangle (13.15,19.6);
\node[pkglabel] at (8.15,19.6) {Ensalamento};
\draw[pkg] (8.15,9.15) rectangle (13.15,12.25);
\node[pkglabel] at (8.15,12.25) {Cadastros};
\draw[pkg] (8.15,5.25) rectangle (13.15,8.85);
\node[pkglabel] at (8.15,8.85) {Administração do sistema};
\draw[pkg] (8.15,-0.25) rectangle (13.15,4.95);
\node[pkglabel] at (8.15,4.95) {Gestão};
\node[uc] (consultar) at (5.35,18.45) {Consultar calendário\\público};
\node[uc] (exportar) at (5.35,16.9) {Exportar calendário\\(.ics)};
\node[uc] (acompanhar) at (5.35,14.55) {Acompanhar e\\cancelar solicitações};
\node[uc] (solicitar) at (5.35,13.25) {Solicitar reserva};
\node[uc] (verificar) at (5.35,11.6) {Verificar\\disponibilidade};
\node[uc] (analisar) at (5.35,9.95) {Analisar solicitação};
\node[uc] (sugerir) at (5.35,8.55) {Sugerir alternativa};
\node[uc] (emnome) at (5.35,7.25) {Solicitar em nome\\de outro usuário};
\node[uc] (receber) at (5.35,4.9) {Receber\\notificações};
\node[uc] (ajuda) at (5.35,3.35) {Consultar central\\de ajuda};
\node[uc] (importar) at (10.65,18.45) {Importar planilha\\de horários};
\node[uc] (resolver) at (10.65,16.55) {Resolver espaços\\da planilha};
\node[uc] (aulas) at (10.65,15.0) {Manter aulas\\recorrentes};
\node[uc] (ocorrencia) at (10.65,13.45) {Cancelar ocorrência\\de aula};
\node[uc] (espacos) at (10.65,11.35) {Manter espaços};
\node[uc] (cursos) at (10.65,10.05) {Manter cursos,\\disciplinas e turmas};
\node[uc] (perfis) at (10.65,7.7) {Gerenciar perfis\\e permissões};
\node[uc] (email) at (10.65,6.15) {Configurar servidor\\de e-mail};
\node[uc] (usuarios) at (10.65,4.05) {Gerenciar usuários};
\node[uc] (painel) at (10.65,2.75) {Consultar painel\\e relatórios};
\node[uc] (auditoria) at (10.65,1.45) {Consultar auditoria};
\node[uc] (comunicado) at (10.65,0.4) {Publicar comunicado};
\draw[rel] (solicitar) -- node[rellabel] {«include»} (verificar);
\draw[rel] (analisar) -- node[rellabel] {«include»} (verificar);
\draw[rel] (sugerir) -- node[rellabel] {«extend»} (analisar);
\draw[rel] (importar) -- node[rellabel] {«include»} (resolver);
\pic at (1.2,17.675) {actor};
\node[hit] (Avis) at (1.2,17.725) {};
\node[lbl] (Lvis) at (1.2,17.125) {Visitante};
\draw (Avis) -- (consultar.west);
\draw (Avis) -- (exportar.west);
\pic at (0.6,13.9) {actor};
\node[hit] (Asol) at (0.6,13.950) {};
\node[lbl] (Lsol) at (0.6,13.350) {Solicitante};
\draw (Asol) -- (acompanhar.west);
\draw (Asol) -- (solicitar.west);
\pic at (1.6,7.5) {actor};
\node[hit] (Aana) at (1.6,7.550) {};
\node[lbl] (Lana) at (1.6,6.950) {Analista de\\reservas};
\draw (Aana) -- (analisar.west);
\draw (Aana) -- (emnome.west);
\pic at (1.2,4.125) {actor};
\node[hit] (Aaut) at (1.2,4.175) {};
\node[lbl] (Laut) at (1.2,3.575) {Usuário\\autenticado};
\draw (Aaut) -- (receber.west);
\draw (Aaut) -- (ajuda.west);
\pic at (15.0,10.65) {actor};
\node[hit] (Ages) at (15.0,10.700) {};
\node[lbl] (Lges) at (15.0,10.100) {Gestor do\\ensalamento};
\draw (Ages) -- (importar.east);
\draw (Ages) -- (aulas.east);
\draw (Ages) -- (ocorrencia.east);
\draw (Ages) -- (espacos.east);
\draw (Ages) -- (cursos.east);
\pic at (15.0,6.25) {actor};
\node[hit] (Asup) at (15.0,6.300) {};
\node[lbl] (Lsup) at (15.0,5.700) {Superadministrador};
\draw (Asup) -- (perfis.east);
\draw (Asup) -- (email.east);
\pic at (15.0,0.6000000000000001) {actor};
\node[hit] (Aadm) at (15.0,0.650) {};
\node[lbl] (Ladm) at (15.0,0.050) {Administrador};
\draw (Aadm) -- (usuarios.east);
\draw (Aadm) -- (painel.east);
\draw (Aadm) -- (auditoria.east);
\draw (Aadm) -- (comunicado.east);
\draw[gen] (Lsol.south) -- (0.6,5.34) -- (1.08,4.88);
\draw[gen] (Lana.south) -- (1.3199999999999998,4.88);
\draw[gen] (Lsup.south) -- (15.0,1.35);
\draw (Ages.east) -- (16.55,10.70) -- (16.55,-1.10);
\draw (Ladm.south) -- (15.0,-1.10);
\draw[gen] (16.55,-1.10) -- (1.2,-1.10) -- (Laut.south);
\end{tikzpicture}
}
    \caption*{Fonte: Elaborado pelo autor (2026).}
    \label{fig:casos_uso}
\end{figure}

O diagrama de classes (Figura~\ref{fig:classes}) apresenta as 18 classes do domínio, com seus atributos principais e as multiplicidades das associações. Classes de apoio, como importação, auditoria e fila de e-mails, foram omitidas para preservar a legibilidade.

O diagrama de sequência (Figura~\ref{fig:sequencia}) detalha a solicitação e a aprovação de uma reserva. Os fragmentos \textit{alt} mostram as duas verificações de conflito de horário: a primeira na solicitação e a segunda na aprovação, quando a restrição do banco de dados impede agendamentos sobrepostos (RNF08).

\begin{figure}[p]
    \centering
    \caption{Diagrama de classes do sistema}
    \resizebox{\textwidth}{!}{% Diagrama de classes do domínio do sistema de agendamento (página em pé).
% Fonte: models.py das apps accounts, schedule, bookings e notices.
% Classes de apoio (importação, auditoria, fila de e-mails) foram omitidas.
\begin{tikzpicture}[
    font=\scriptsize,
    cls/.style={draw, rectangle split, rectangle split parts=3,
                text width=3.4cm, align=left, anchor=north, inner sep=2.5pt,
                rectangle split part fill={black!8,white,white}},
    cls2/.style={draw, rectangle split, rectangle split parts=2,
                text width=3.4cm, align=left, anchor=north, inner sep=2.5pt,
                rectangle split part fill={black!8,white}},
    m/.style={font=\tiny, inner sep=1.2pt},
    rl/.style={font=\tiny\itshape, inner sep=1.2pt, fill=white},
    comp/.style={{Diamond[length=2.6mm, width=1.8mm]}-},
]
\newcommand{\cname}[1]{\centering\textbf{\footnotesize #1}\par}
\def\ca{1.85}\def\cb{6.85}\def\cc{11.45}\def\cd{16.05}
\def\ra{21.0}\def\rb{18.4}\def\rc{15.2}\def\rd{9.8}\def\re{6.9}

% ---------------- coluna 1: acesso ----------------
\node[cls2] (permissao) at (\ca,\ra) {\cname{Permissão}
  \nodepart{second} - codigo: texto\\- descricao: texto};
\node[cls2] (perfil) at (\ca,\rb) {\cname{Perfil}
  \nodepart{second} - codigo: texto\\- nome: texto\\- descricao: texto};
\node[cls] (usuario) at (\ca,\rc) {\cname{Usuário}
  \nodepart{second} - nome: texto\\- email: texto\\- nome\_usuario: texto\\- ativo: lógico\\- superadmin: lógico
  \nodepart{third} + tem\_permissao(cod)};
\node[cls2] (notificacao) at (\ca,\rd) {\cname{Notificação}
  \nodepart{second} - titulo: texto\\- mensagem: texto\\- lida\_em: data-hora};

% ---------------- coluna 2: agenda ----------------
\node[cls2] (outronome) at (\cb,\ra) {\cname{Outro nome}
  \nodepart{second} - nome: texto};
\node[cls2] (agendamento) at (\cb,\rb) {\cname{Agendamento}
  \nodepart{second} - inicio: data-hora\\- fim: data-hora\\- descricao: texto\\- qtd\_alunos: inteiro};
\node[cls] (solicitacao) at (\cb,\rc) {\cname{Solicitação}
  \nodepart{second} - data: data\\- inicio: hora\\- fim: hora\\- descricao: texto\\- qtd\_alunos: inteiro\\- situacao: texto
  \nodepart{third} + submeter()\\+ aprovar(analista)\\+ rejeitar(motivo)\\+ sugerir(alternativa)\\+ cancelar(motivo)};
\node[cls2] (vinculo) at (\cb,\rd) {\cname{Vínculo docente}
  \nodepart{second} - atual: lógico};
\node[cls2] (docente) at (\cb,\re) {\cname{Docente}
  \nodepart{second} - nome: texto};

% ---------------- coluna 3: espaços e grade ----------------
\node[cls2] (espaco) at (\cc,\rb) {\cname{Espaço}
  \nodepart{second} - nome: texto\\- origem: texto\\- ativo: lógico};
\node[cls2] (ensalamento) at (\cc,\rc) {\cname{Ensalamento}
  \nodepart{second} - dia\_semana: texto\\- origem: texto\\- vigente\_de: data\\- vigente\_ate: data};
\node[cls2] (disciplina) at (\cc,\rd) {\cname{Disciplina}
  \nodepart{second} - nome: texto};
\node[cls2] (oferta) at (\cc,\re) {\cname{Oferta de disciplina}
  \nodepart{second} \vphantom{x}};

% ---------------- coluna 4 ----------------
\node[cls2] (faixa) at (\cd,\ra) {\cname{Faixa de horário}
  \nodepart{second} - inicio: hora\\- fim: hora};
\node[cls2] (semestre) at (\cd,\rb) {\cname{Semestre}
  \nodepart{second} - codigo: texto\\- inicio: data\\- fim: data};
\node[cls2] (ocorrencia) at (\cd,\rc) {\cname{Ocorrência cancelada}
  \nodepart{second} - data: data\\- tipo: texto};
\node[cls2] (turma) at (\cd,\rd) {\cname{Turma}
  \nodepart{second} - rotulo: texto};
\node[cls2] (curso) at (\cd,\re) {\cname{Curso}
  \nodepart{second} - nome: texto\\- etapa: texto};

% ================= associações =================
% --- acesso (verticais na coluna 1)
\draw (perfil.north) -- node[rl, right] {concede} (permissao.south)
      node[m, pos=0.12, left] {*} node[m, pos=0.88, left] {*};
\draw (usuario.north) -- node[rl, right] {tem} (perfil.south)
      node[m, pos=0.12, left] {*} node[m, pos=0.88, left] {0..1};
\draw (usuario.south) -- node[rl, right] {recebe} (notificacao.north)
      node[m, pos=0.08, left] {1} node[m, pos=0.92, left] {*};

% --- usuário e solicitação (horizontais)
\path (usuario.north east) ++(0,-0.75) coordinate (s1);
\path (usuario.north east) ++(0,-1.75) coordinate (s2);
\draw (s1) -- node[rl, above] {solicita} (solicitacao.west |- s1)
      node[m, pos=0.1, below] {1} node[m, pos=0.9, below] {*};
\draw (s2) -- node[rl, above] {analisa} (solicitacao.west |- s2)
      node[m, pos=0.12, below] {0..1} node[m, pos=0.9, below] {*};

% --- agenda (verticais na coluna 2)
\draw (solicitacao.north) -- node[rl, right] {gera} (agendamento.south)
      node[m, pos=0.12, left] {1} node[m, pos=0.88, left] {0..1};
\draw (solicitacao.south) -- node[rl, right] {justifica} (vinculo.north)
      node[m, pos=0.08, left] {*} node[m, pos=0.92, left] {0..1};
\draw (docente.north) -- node[rl, right] {assume} (vinculo.south)
      node[m, pos=0.12, left] {1} node[m, pos=0.88, left] {*};

% --- espaço
\path (agendamento.north east) ++(0,-0.9) coordinate (a1);
\draw (a1) -- node[rl, above] {ocupa} (espaco.west |- a1)
      node[m, pos=0.1, below] {*} node[m, pos=0.9, below] {1};
\draw (solicitacao.north east) -- node[rl, sloped, below] {é pedida em} (espaco.south west)
      node[m, pos=0.1, above left] {*} node[m, pos=0.85, above left] {1};
\draw[comp] (espaco.north west) -- (outronome.south east)
      node[m, pos=0.85, below right] {*};

% --- ensalamento
\draw (ensalamento.north) -- node[rl, right] {usa} (espaco.south)
      node[m, pos=0.12, left] {*} node[m, pos=0.88, left] {0..1};
\draw (ensalamento.south) -- (disciplina.north)
      node[m, pos=0.06, left] {*} node[m, pos=0.94, left] {1};
\path (ensalamento.north east) ++(0,-0.9) coordinate (e1);
\draw[comp] (e1) -- (ocorrencia.west |- e1)
      node[m, pos=0.85, above] {*};
\draw ([yshift=-3mm]ensalamento.north east) -- (semestre.south west)
      node[m, pos=0.15, below right] {*} node[m, pos=0.85, below right] {1};
\draw (ensalamento.north east) -- (faixa.south west)
      node[m, pos=0.06, left] {*} node[m, pos=0.95, left] {1};
\draw (ensalamento.south east) -- (turma.north west)
      node[m, pos=0.12, below left] {*} node[m, pos=0.88, below left] {1};
\draw (docente.east) -- (\cb+2.3,\re-0.5) -- node[rl, sloped, above, pos=0.55] {ministra} (\cb+2.3,\rc-1.75)
      -- (ensalamento.south west |- 0,\rc-1.75);
\node[m, above right] at (docente.east) {*};
\node[m, above left] at (ensalamento.south west |- 0,\rc-1.75) {*};

% --- grade
\draw (semestre.north) -- node[rl, right] {define} (faixa.south)
      node[m, pos=0.12, left] {1} node[m, pos=0.88, left] {*};
\draw (vinculo.south east) -- (oferta.north west)
      node[m, pos=0.12, above right] {*} node[m, pos=0.85, above right] {1};
\draw (oferta.north) -- (disciplina.south)
      node[m, pos=0.12, left] {*} node[m, pos=0.88, left] {1};
\draw (oferta.north east) -- (turma.south west)
      node[m, pos=0.12, above left] {*} node[m, pos=0.88, above left] {1};
\draw (curso.north) -- node[rl, right] {possui} (turma.south)
      node[m, pos=0.12, left] {1} node[m, pos=0.88, left] {*};
\draw (oferta.south) -- (\cc,\re-1.6) -- (\cd+2.3,\re-1.6) -- (\cd+2.3,\rb-0.8) -- (semestre.east |- 0,\rb-0.8)
      node[m, pos=0.9, above] {1};
\node[m, below left] at (oferta.south) {*};
\end{tikzpicture}
}
    \caption*{Fonte: Elaborado pelo autor (2026).}
    \label{fig:classes}
\end{figure}

\begin{figure}[p]
    \centering
    \caption{Diagrama de sequência da solicitação e aprovação de reserva}
    \resizebox{\textwidth}{!}{% Diagrama de sequência: solicitar e aprovar uma reserva.
% Fonte: bookings/models.py (submit, approve, Booking.objects.book) e notices/events.py (announce).
\begin{tikzpicture}[
    font=\scriptsize,
    obj/.style={draw, fill=black!8, minimum height=0.6cm, inner sep=3pt, font=\footnotesize},
    life/.style={dashed, black!55},
    call/.style={-{Latex[length=2mm]}},
    ret/.style={dashed, -{Stealth[open, length=2mm]}},
    async/.style={-{Stealth[open, length=2mm]}},
    lab/.style={above, inner sep=1.5pt, fill=white, fill opacity=0.9, text opacity=1},
    frag/.style={draw, black!70},
    guard/.style={font=\scriptsize\itshape, anchor=north west, inner sep=2pt},
    actor/.pic={
        \draw (0,0.55) circle (0.17);
        \draw (0,0.38) -- (0,-0.1) (-0.3,0.22) -- (0.3,0.22)
              (0,-0.1) -- (-0.25,-0.5) (0,-0.1) -- (0.25,-0.5);
    },
]
\def\xs{0.8}\def\xa{2.8}\def\xi{5.0}\def\xq{7.6}\def\xo{10.2}\def\xg{12.6}\def\xn{15.0}
\def\ytop{0}\def\ybot{-17.6}

% ---------- quadro sd ----------
\draw[frag] (-0.2,2.0) rectangle (16.2,\ybot-0.3);
\node[anchor=north west, inner sep=2pt, font=\scriptsize\bfseries] (sd) at (-0.2,2.0) {sd Solicitar e aprovar reserva};
\draw[frag] (sd.south west) -| ([xshift=2mm]sd.east) -- ([xshift=2mm]sd.north east);

% ---------- participantes ----------
\pic at (\xs,0.55) {actor}; \node at (\xs,-0.15) {Solicitante};
\pic at (\xa,0.55) {actor}; \node at (\xa,-0.15) {Analista};
\node[obj] at (\xi,0.5) {Interface web};
\node[obj] at (\xq,0.5) {:Solicitação};
\node[obj] at (\xo,0.5) {:Ocupação};
\node[obj] at (\xg,0.5) {:Agendamento};
\node[obj] at (\xn,0.5) {:Notificações};
\foreach \x/\y in {\xs/-0.4,\xa/-0.4,\xi/0.2,\xq/0.2,\xo/0.2,\xg/0.2,\xn/0.2}
    \draw[life] (\x,\y) -- (\x,\ybot);

% ================= 1. solicitação =================
\draw[call] (\xs,-1.0) -- node[lab] {solicitar(espaço, data, horário)} (\xi,-1.0);
\draw[call] (\xi,-1.6) -- node[lab] {submeter(dados)} (\xq,-1.6);
\draw[call] (\xq,-2.1) -- ++(0.6,0) -- ++(0,-0.4) node[right, midway, inner sep=2pt] {verificar permissão} -- (\xq+0.08,-2.5);
\draw[call] (\xq,-3.1) -- node[lab] {verificar\_conflito(espaço, horário)} (\xo,-3.1);
\draw[ret]  (\xo,-3.7) -- node[lab] {resultado} (\xq,-3.7);

% alt: conflito na solicitação
\draw[frag] (0.1,-4.0) rectangle (15.9,-8.7);
\node[draw, frag, fill=white, anchor=north west, inner sep=2pt, font=\scriptsize\bfseries] at (0.1,-4.0) {alt};
\node[guard] at (1.0,-4.0) {[horário ocupado]};
\draw[ret] (\xq,-4.6) -- node[lab] {erro: horário ocupado} (\xi,-4.6);
\draw[ret] (\xi,-5.1) -- node[lab] {mensagem de erro} (\xs,-5.1);
\draw[frag, dashed] (0.1,-5.4) -- (15.9,-5.4);
\node[guard] at (1.0,-5.4) {[horário livre]};
\draw[call] (\xq,-5.9) -- ++(0.6,0) -- ++(0,-0.4) node[right, midway, inner sep=2pt] {gravar como PENDENTE e auditar} -- (\xq+0.08,-6.3);
\draw[call] (\xq,-6.9) -- node[lab] {anunciar(SOLICITADA)} (\xn,-6.9);
\draw[async] (\xn,-7.5) -- node[lab, pos=0.75] {aviso em tempo real} (\xa,-7.5);
\draw[ret] (\xq,-8.0) -- node[lab] {solicitação} (\xi,-8.0);
\draw[ret] (\xi,-8.4) -- node[lab] {enviada} (\xs,-8.4);

% ================= 2. análise =================
\draw[call] (\xa,-9.3) -- node[lab] {aprovar(solicitação)} (\xi,-9.3);
\draw[call] (\xi,-9.9) -- node[lab] {aprovar(analista)} (\xq,-9.9);
\draw[call] (\xq,-10.4) -- ++(0.6,0) -- ++(0,-0.4) node[right, midway, inner sep=2pt] {verificar permissão e bloquear registro} -- (\xq+0.08,-10.8);
\draw[call] (\xq,-11.4) -- node[lab] {verificar\_conflito(espaço, horário)} (\xo,-11.4);
\draw[ret]  (\xo,-12.0) -- node[lab] {livre} (\xq,-12.0);
\draw[call] (\xq,-12.6) -- node[lab] {criar(espaço, início, fim)} (\xg,-12.6);

% alt: restrição do banco
\draw[frag] (0.1,-12.9) rectangle (15.9,-17.3);
\node[draw, frag, fill=white, anchor=north west, inner sep=2pt, font=\scriptsize\bfseries] at (0.1,-12.9) {alt};
\node[guard] at (1.0,-12.9) {[restrição de sobreposição violada]};
\draw[ret] (\xg,-13.5) -- node[lab] {horário indisponível} (\xq,-13.5);
\draw[ret] (\xq,-14.0) -- node[lab] {erro} (\xi,-14.0);
\draw[ret] (\xi,-14.5) -- node[lab] {mensagem de erro} (\xa,-14.5);
\draw[frag, dashed] (0.1,-14.8) -- (15.9,-14.8);
\node[guard] at (1.0,-14.8) {[agendamento criado]};
\draw[ret] (\xg,-15.3) -- node[lab] {agendamento} (\xq,-15.3);
\draw[call] (\xq,-15.9) -- node[lab] {anunciar(APROVADA)} (\xn,-15.9);
\draw[async] (\xn,-16.5) -- node[lab, pos=0.85] {aviso em tempo real} (\xs,-16.5);
\draw[ret] (\xi,-17.0) -- node[lab] {aprovada} (\xa,-17.0);
\end{tikzpicture}
}
    \caption*{Fonte: Elaborado pelo autor (2026).}
    \label{fig:sequencia}
\end{figure}

%-----------------------------------
\subsubsection{Implementação}
\label{sec:implementacao}

A quarta etapa foi a implementação. O IFAgenda foi desenvolvido como uma aplicação web, com as tecnologias e a arquitetura descritas a seguir, e constituiu uma versão funcional, executada em ambiente local e distribuída em contêineres.

A implementação utiliza Python e o \textit{framework} Django, cujo padrão arquitetural MTV e funcionamento estão descritos nas Seções~\ref{sec:mvc_mtv} e~\ref{sec:django}. O \textit{framework} fornece o ORM, o sistema de autenticação, o roteamento e o mecanismo de \textit{templates} utilizados no projeto. A separação entre dados, lógica e apresentação também organiza o código e concentra as regras de reserva em componentes próprios.

A Tabela~\ref{tab:funcionalidades} resume as funcionalidades do sistema, agrupadas pelos mesmos módulos dos requisitos funcionais.

\begin{table}[!htb]
	\centering
	\caption{Funcionalidades do sistema por módulo}
	\renewcommand{\arraystretch}{1.25}
	\newcommand{\fitem}[1]{\par\hangindent=1em\hangafter=1\noindent\makebox[1em][l]{\textbullet}#1}
	\begin{tabular}{@{}>{\raggedright\arraybackslash}p{4.0cm}>{\raggedright\arraybackslash}p{9.8cm}@{}}
	\hline
		\textbf{Módulo} & \textbf{Funcionalidades principais} \\ \hline
		Ensalamento & \fitem{Importação da planilha de horários, com prévia e confirmação}\fitem{Associação dos espaços da planilha aos já cadastrados}\fitem{Aulas recorrentes: alteração, encerramento e cancelamento de ocorrências} \\ \arrayrulecolor{black!25}\hline\arrayrulecolor{black}
		Espaços, cursos, disciplinas e turmas & \fitem{Cadastro e manutenção} \\ \arrayrulecolor{black!25}\hline\arrayrulecolor{black}
		Usuários e perfis & \fitem{Cadastro individual ou em lote (CSV)}\fitem{Definição e recuperação de senha}\fitem{Perfis com permissões definidas pelo superadministrador} \\ \arrayrulecolor{black!25}\hline\arrayrulecolor{black}
		Agenda de reservas & \fitem{Solicitação, inclusive em nome de outro usuário}\fitem{Análise: aprovação, rejeição ou sugestão de alternativa}\fitem{Acompanhamento das próprias solicitações} \\ \arrayrulecolor{black!25}\hline\arrayrulecolor{black}
		Notificações e comunicados & \fitem{Avisos em tempo real}\fitem{Comunicados para todos, para um perfil ou para usuários escolhidos} \\ \arrayrulecolor{black!25}\hline\arrayrulecolor{black}
		Consulta pública & \fitem{Calendário por data, com filtros por sala, docente, curso e turma}\fitem{Horário original da planilha}\fitem{Exportação do calendário (.ics)} \\ \arrayrulecolor{black!25}\hline\arrayrulecolor{black}
		Painel, relatórios e auditoria & \fitem{Painel de pendências e solicitações}\fitem{Mapa de ocupação por dia e horário}\fitem{Relatório de solicitações}\fitem{Registro das ações, sem alteração ou exclusão} \\ \arrayrulecolor{black!25}\hline\arrayrulecolor{black}
		Configurações & \fitem{Servidor de envio de e-mails configurado pela interface} \\ \arrayrulecolor{black!25}\hline\arrayrulecolor{black}
		Central de ajuda & \fitem{Tutoriais por tela, filtrados pelas permissões do usuário} \\ \hline
	\end{tabular}
	\label{tab:funcionalidades}
\end{table}

Para a camada de dados, foi utilizado o PostgreSQL como sistema gerenciador de banco de dados relacional. O SQLite foi mantido como alternativa adequada para desenvolvimento local, mas o PostgreSQL atende melhor ao cenário de execução definido, que envolve acessos simultâneos, transações e restrições de integridade para impedir reservas sobrepostas. O Django oferece suporte ao PostgreSQL por meio do ORM e das transações do banco \cite{djangodocs}. No projeto, o acesso aos dados é feito pelo ORM, enquanto as restrições de integridade são mantidas no esquema relacional.

A camada de apresentação foi construída com os \textit{templates} do Django combinados ao Tailwind CSS. Essa escolha agiliza o desenvolvimento da interface, mantém a consistência visual entre as telas e reduz o custo de manutenção, permitindo ajustes visuais sem alteração nas \textit{views} ou nos \textit{models}.

Para a autenticação dos usuários, foi utilizada a biblioteca \textit{django-allauth}, escolhida por fornecer prontos os fluxos de registro, login, recuperação de senha e gerenciamento de sessões \cite{pennersr2024}. O controle de acesso descrito na Seção~\ref{sec:rbac} foi implementado com perfis configuráveis. Cada permissão corresponde a uma ação sobre um recurso do sistema, como solicitar uma reserva ou editar um espaço, e o sistema define 47 permissões, distribuídas em 13 módulos. Um perfil é um conjunto dessas permissões, montado pelo superadministrador. Os perfis Administrador e Docente vêm na carga inicial apenas como exemplo e podem ser alterados. Outras categorias, como os técnicos administrativos, recebem um perfil com as permissões de que precisam. A verificação ocorre sempre no servidor, e a consulta pública dispensa autenticação.

As notificações automáticas são exibidas no próprio sistema, em um ícone no cabeçalho que se atualiza em tempo real, sem recarregar a página (RNF11). Ao receber uma solicitação de reserva, o sistema notifica os responsáveis pela análise; após a decisão, o solicitante recebe o resultado (RF09 e RF10). O e-mail é usado apenas para a definição e a recuperação de senha, com o servidor de envio configurado pela interface (RF27) e as mensagens enviadas em segundo plano (RNF10) \cite{djangodocs}.

As notificações em tempo real foram implementadas com o Django Channels, extensão que permite ao Django manter conexões WebSocket com o navegador \cite{channelsdocs}. Tarefas demoradas, como a importação da planilha e o envio de e-mails, são executadas por um processo próprio em segundo plano. Esse processo lê uma fila registrada no banco e repete o envio em caso de falha. A aplicação, o banco de dados e o processo de tarefas são distribuídos em contêineres orquestrados pelo Docker Compose \cite{dockerdocs}.

O presente trabalho limita-se ao desenvolvimento de uma versão funcional do sistema para execução em ambiente local, acessível por navegador. A hospedagem e a implantação em infraestrutura institucional não fazem parte do escopo.
%-----------------------------------
\subsubsection{Validação}
\label{sec:validacao}

A validação realizada neste trabalho ocorreu em duas frentes. A primeira correspondeu à validação do alinhamento da solução ao processo institucional, conduzida com a Diretora de Ensino, principal interessada e responsável pelos agendamentos. A segunda foi exclusivamente técnica e consistiu em testes automatizados, unitários e de integração, conforme os conceitos da Seção~\ref{sec:testes}. Os testes unitários verificaram regras e componentes de forma isolada. Os testes de integração verificaram o comportamento conjunto entre as camadas da aplicação, o banco de dados e os fluxos acessados pelos usuários.

Os testes são escritos com o \textit{framework} de testes do Django. Cada caso é executado sobre um banco de dados temporário, separado dos dados utilizados durante o desenvolvimento \cite{djangodocs}. Esse procedimento permite preparar os dados necessários para cada cenário e verificar as alterações realizadas sem modificar a base utilizada pela aplicação. A versão desenvolvida conta com mais de 1.700 casos de teste, distribuídos entre os módulos e regras de negócio do sistema.

Entre as situações verificadas estiveram o bloqueio de reservas sobrepostas, inclusive quando solicitações foram analisadas simultaneamente; a aplicação das permissões conforme o perfil do usuário; a importação e a reimportação da planilha de horários; as alterações em aulas recorrentes; o envio de notificações; a consulta pública sem exposição de dados pessoais; e a imutabilidade dos registros de auditoria. A Tabela~\ref{tab:testes} relacionou essas verificações aos requisitos correspondentes. Para a validação técnica, foi adotado como critério a execução de todos os testes automatizados sem falhas. Não foram realizados testes formais de usabilidade com usuários finais. Portanto, os testes automatizados avaliaram o funcionamento técnico do sistema, enquanto a adequação aos requisitos institucionais foi analisada nas reuniões com a Diretora de Ensino.

\begin{table}[!htb]
	\centering
	\caption{Regras de negócio verificadas pelos testes}
	\begin{tabular}{@{}>{\raggedright\arraybackslash}p{9.8cm}>{\raggedright\arraybackslash}p{4.0cm}@{}}
	\hline
		\textbf{Regra verificada} & \textbf{Requisitos} \\ \hline
		Bloqueio de reservas sobrepostas, inclusive em análises simultâneas & RF06, RF07, RNF08 \\
		Acesso a cada funcionalidade conforme as permissões do perfil & RF18, RNF01 \\
		Importação e reimportação da planilha de horários & RF01, RF02 \\
		Alteração, encerramento e cancelamento de aulas recorrentes & RF14--RF16 \\
		Envio das notificações ao solicitante e aos responsáveis & RF09, RF10 \\
		Consulta pública sem exposição de dados pessoais & RF11 \\
		Imutabilidade dos registros de auditoria & RF26, RNF09 \\ \hline
	\end{tabular}
	\label{tab:testes}
\end{table}

%---------------------------------
% SEÇÃO DE CRONOGRAMA
%---------------------------------
\needspace{9\baselineskip}
